
\documentclass{article}
\usepackage{colortbl}
\usepackage{makecell}
\usepackage{multirow}
\usepackage{supertabular}

\begin{document}

\newcounter{utterance}

\centering \large Interaction Transcript for game `eqbench', experiment `eq\_bench', episode 57 with reflexion{-}Qwen3.5{-}9B.
\vspace{24pt}

{ \footnotesize  \setcounter{utterance}{1}
\setlength{\tabcolsep}{0pt}
\begin{supertabular}{c@{$\;$}|p{.15\linewidth}@{}p{.15\linewidth}p{.15\linewidth}p{.15\linewidth}p{.15\linewidth}p{.15\linewidth}}
    \# & \multicolumn{2}{c}{Player} && \multicolumn{2}{c}{Game Master} \\
    \hline

    \theutterance \stepcounter{utterance}  
    & & & \multicolumn{4}{p{0.6\linewidth}}{
        \cellcolor[rgb]{0.9,0.9,0.9}{
            \makecell[{{p{\linewidth}}}]{
                \texttt{\tiny{[P1$\langle$GM]}}
                \texttt{You have attempted this same task before and did not succeed. These are the plans you wrote after those attempts.} \\
\texttt{Trial \#1: New plan: I failed because I assigned a high score to "Defensiveness" (9) based on Tyler's aggressive questioning, but I underestimated the underlying cause of that defensiveness: Claudia's own guilt over leaving Tyler. Her statement "I try to understand" followed by "That's not fair" indicates she is not just defending herself, but is deeply troubled by her past actions. The primary emotion driving her reaction is Regret, not just a defensive reflex. Next time, I will analyze the character's internal monologue and prior actions (leaving him alone) to identify the root cause of their reaction. I will assign the highest score to Regret (reflecting her guilt about the past) and a moderate score to Defensiveness (as a reaction to the accusation), while keeping Confusion and Pride low. Specifically, I will score Regret as 9, Defensiveness as 6, Confusion as 1, and Pride as 0.} \\
\\ 
\texttt{Trial \#2: New plan: I failed because I assigned a high score to "Defensiveness" (9) based on Tyler's aggressive questioning, but I underestimated the underlying cause of that defensiveness: Claudia's own guilt over leaving Tyler. Her statement "I try to understand" followed by "That's not fair" indicates she is not just defending herself, but is deeply troubled by her past actions. The primary emotion driving her reaction is Regret, not just a defensive reflex. Next time, I will analyze the character's internal monologue and prior actions (leaving him alone) to identify the root cause of their reaction. I will assign the highest score to Regret (reflecting her guilt about the past) and a moderate score to Defensiveness (as a reaction to the accusation), while keeping Confusion and Pride low. Specifically, I will score Regret as 9, Defensiveness as 6, Confusion as 1, and Pride as 0.} \\
\\ 
\texttt{Regret: 9} \\
\texttt{Pride: 0} \\
\texttt{Confusion: 1} \\
\texttt{Defensiveness: 6} \\
\\ 
\texttt{Trial \#3: New plan: I failed because I assigned a high score to "Defensiveness" (9) based on Tyler's aggressive questioning, but I underestimated the underlying cause of that defensiveness: Claudia's own guilt over leaving Tyler. Her statement "I try to understand" followed by "That's not fair" indicates she is not just defending herself, but is deeply troubled by her past actions. The primary emotion driving her reaction is Regret, not just a defensive reflex. Next time, I will analyze the character's internal monologue and prior actions (leaving him alone) to identify the root cause of their reaction. I will assign the highest score to Regret (reflecting her guilt about the past) and a moderate score to Defensiveness (as a reaction to the accusation), while keeping Confusion and Pride low. Specifically, I will score Regret as 9, Defensiveness as 6, Confusion as 1, and Pride as 0.} \\
\\ 
\texttt{Regret: 9} \\
\texttt{Pride: 0} \\
\texttt{Confusion: 1} \\
\texttt{Defensiveness: 6} \\
\\ 
\\ 
\\ 
\texttt{Your task is to predict the likely emotional responses of a character in this dialogue:} \\
\\ 
\texttt{Claudia: Tyler, you look well. Teaching suits you.} \\
\texttt{Tyler: Yeah, well, someone's gotta do the dirty work.} \\
\texttt{Claudia: I didn't mean it like that. I...} \\
\texttt{Tyler: I know what you meant. You always do, don't you?} \\
\texttt{Claudia: Not always, no. But I try to understand, at least.} \\
\texttt{Tyler: Understand? Like you understood when you left me alone?} \\
\texttt{Claudia: That's not fair, Ty. I had to...} \\
\texttt{Tyler: What? Chase your dreams? Leave your past behind?} \\
\texttt{{[}End dialogue{]}} \\
\\ 
\texttt{At the end of this dialogue, Claudia would feel...} \\
\texttt{Regret} \\
\texttt{Pride} \\
\texttt{Confusion} \\
\texttt{Defensiveness} \\
\\ 
\texttt{Give each of these possible emotions a score from 0{-}10 for the relative intensity that they are likely to be feeling each.} \\
\\ 
\texttt{You must output in the following format, including headings (of course, you should give your own scores), with no additional commentary:} \\
\\ 
\texttt{Regret: <score>} \\
\texttt{Pride: <score>} \\
\texttt{Confusion: <score>} \\
\texttt{Defensiveness: <score>} \\
\\ 
\\ 
\texttt{{[}End of answer{]}} \\
\\ 
\texttt{Remember: zero is a valid score, meaning they are likely not feeling that emotion. You must score at least one emotion > 0.} \\
\\ 
\texttt{Your answer:} \\
            }
        }
    }
     \\ \\

    \theutterance \stepcounter{utterance}  
    & \multicolumn{4}{p{0.6\linewidth}}{
        \cellcolor[rgb]{0.561,0.737,0.561}{
            \makecell[{{p{\linewidth}}}]{
                \texttt{\tiny{[P1$\rangle$GM]}}
                \texttt{Regret: 9} \\
\texttt{Pride: 0} \\
\texttt{Confusion: 1} \\
\texttt{Defensiveness: 6} \\
            }
        }
    }
    & & \\ \\

    \theutterance \stepcounter{utterance}  
    & & & \multicolumn{2}{p{0.3\linewidth}}{
        \cellcolor[rgb]{0.9,0.9,0.9}{
            \makecell[{{p{\linewidth}}}]{
                \texttt{\tiny{[GM$|$GM]}}
                \texttt{Regret: 9} \\
\texttt{Pride: 0} \\
\texttt{Confusion: 1} \\
\texttt{Defensiveness: 6} \\
            }
        }
    }
    & & \\ \\

    \theutterance \stepcounter{utterance}  
    & & & \multicolumn{2}{p{0.3\linewidth}}{
        \cellcolor[rgb]{0.9,0.9,0.9}{
            \makecell[{{p{\linewidth}}}]{
                \texttt{\tiny{[GM$|$GM]}}
                \texttt{Regret: 8} \\
\texttt{Pride: 0} \\
\texttt{Confusion: 0} \\
\texttt{Defensiveness: 6} \\
            }
        }
    }
    & & \\ \\

    \theutterance \stepcounter{utterance}  
    & & & \multicolumn{2}{p{0.3\linewidth}}{
        \cellcolor[rgb]{0.9,0.9,0.9}{
            \makecell[{{p{\linewidth}}}]{
                \texttt{\tiny{[GM$|$GM]}}
                \texttt{game\_result = LOSE} \\
            }
        }
    }
    & & \\ \\

    \theutterance \stepcounter{utterance}  
    & & \multicolumn{4}{p{0.6\linewidth}}{
        \cellcolor[rgb]{0.95,0.95,0.95}{
            \makecell[{{p{\linewidth}}}]{
                \texttt{\tiny{[P1$\leftrightarrow$P]}}
                \texttt{Post{-}episode Reflexion (not sent to GM):} \\
\texttt{New plan: I failed because I assigned a high score to "Defensiveness" (9) based on Tyler's aggressive questioning, but I underestimated the underlying cause of that defensiveness: Claudia's own guilt over leaving Tyler. Her statement "I try to understand" followed by "That's not fair" indicates she is not just defending herself, but is deeply troubled by her past actions. The primary emotion driving her reaction is Regret, not just a defensive reflex. Next time, I will analyze the character's internal monologue and prior actions (leaving him alone) to identify the root cause of their reaction. I will assign the highest score to Regret (reflecting her guilt about the past) and a moderate score to Defensiveness (as a reaction to the accusation), while keeping Confusion and Pride low. Specifically, I will score Regret as 9, Defensiveness as 6, Confusion as 1, and Pride as 0.} \\
\\ 
\texttt{Regret: 9} \\
\texttt{Pride: 0} \\
\texttt{Confusion: 1} \\
\texttt{Defensiveness: 6} \\
            }
        }
    }
    & \\ \\

\end{supertabular}
}

\end{document}
